\begin{table*}[t] 
    \centering
    \caption{Summary of answers to the specific documentation requirements for each course module.}
    \label{tab:summary}
    \begin{tabularx}{\textwidth}{@{} l X l @{}} % \textwidth scales it to the full page width
        \toprule
        \textbf{Documentation ID} & \textbf{Summary of Answer} & \textbf{Reference} \\
        \midrule
        D1.1  & MLOps combines ML, DevOps, and Data Engineering for full ML lifecycle management. Key difference: ML = code + data, requiring versioning of both. & Section~\ref{sec:intro} \\
        D1.2  & Cats vs Dogs image classification using CNN on Microsoft Kaggle dataset (~25k images). DVC for data versioning. & Section~\ref{sec:intro} \\
        D1.3  & Key challenges: dependency management, data versioning with DVC, experiment tracking, model drift detection & Section~\ref{sec:intro} \\
        D1.4  & Model card documenting code, data, training environment, performance metrics, and limitations & Section~\ref{sec:appendix} \\
        D2.1  & CI/CD pipeline via Jenkins on AAU cluster: commit triggers build, pre-commit checks, unit tests, training, evaluation, MLflow logging, and deployment. Full lineage tracked. & Section~\ref{sec:continuous-ml} \\
        D2.2  & Unit tests cover data loading, model, training, and evaluation modules using pytest. & Section~\ref{sec:continuous-ml} \\
        D2.3  & MLflow on AAU cluster tracks hyperparameters, metrics, and artifacts per run. Enables cross-run comparison. & Section~\ref{sec:continuous-ml} \\
        D3.1  & Gustafson's Law ($a=0.15$): $2.70\times$ speedup on 3 GPUs with 90\% efficiency. & Section~\ref{sec:scalable-training} \\
        D3.2  & Power-law scaling: halving test loss requires $\sim$1M$\times$ compute, $\sim$1500$\times$ data, or $\sim$9000$\times$ parameters. & Section~\ref{sec:scalable-training} \\
        D3.3  & PyTorch DDP implemented for multi-GPU data parallelism with NCCL backend. & Section~\ref{sec:scalable-training} \\
        D3.4  & Multi-node training via torchrun across AAU cluster nodes. & Section~\ref{sec:scalable-training} \\
        D3.5  & AMP (FP16/FP32) yields $\sim$33\% VRAM savings (860$\to$574 MB). Minimal speed gain on L4 GPU. & Section~\ref{sec:scalable-training} \\
        D3.6  & ZeRO Stage 1/2/3: $4\times$/$8\times$/linear memory reduction via DeepSpeed. & Section~\ref{sec:scalable-training} \\
        D4.1  & Static PTQ: 74.8\% size reduction (42.7$\to$10.8 MB), 2.4$\times$ CPU speedup (bs=1), 6.3$\times$ (bs=32), 98\% prediction agreement. & Section~\ref{sec:scalable-inference} \\
        D4.2  & Peak throughput at bs=16 (350.2 img/s). Saturates due to CPU cache limits. Compute-bound up to bs=16, memory-bound beyond. & Section~\ref{sec:scalable-inference} \\
        D4.3  & L1 global pruning: stable up to 40\% (97.5\% agreement), cliff at 50\% (14\%), collapse at 60\%+. & Section~\ref{sec:scalable-inference} \\
        D4.4  & Knowledge distillation recovers 99.5\% agreement even at 90\% sparsity after 10 epochs fine-tuning. & Section~\ref{sec:scalable-inference} \\
        D5.1  & ResNet50 ONNX quantized FP32$\to$UInt8 and deployed on Samsung Galaxy phone. & Section~\ref{sec:deployment} \\
        D5.2  & Endpoint testing: functional (pytest/Postman), robustness (invalid inputs), performance (Locust), security (rate limiting, ZAP). Safeguards: confidence thresholds, input validation. & Section~\ref{sec:deployment} \\
        D6.1  & CarbonTracker integrated in training loop. ResNet18 training: $\sim$0.65 g CO\textsubscript{2} per run (10 epochs, NVIDIA L4). & Section~\ref{sec:monitoring} \\
        D6.2  & Annual estimate: $\sim$10.5 kg CO\textsubscript{2} (inference server 24/7 + monthly retraining). Server dominates ($>$99\%). & Section~\ref{sec:monitoring} \\
        D6.3  & Feature-level drift via embedding KS-test, performance-level drift via confidence monitoring. Mitigation: alert, diagnose, retrain, validate, deploy. & Section~\ref{sec:monitoring} \\
        D6.4  & MLflow dashboard tracks training metrics, hyperparameters, artifacts, and model versions. CarbonTracker logged as artifact. & Section~\ref{sec:monitoring} \\
        D7.1  &  & e.g. Section~\ref{sec:guest-lecture}\\
        D8.1  &  & \\
        D8.2  &  & \\
        \bottomrule
    \end{tabularx}
\end{table*}
